\documentclass[a4paper,UTF8,12pt]{ctexart}
\usepackage{Mystyle}
\usepackage{translation-adapter}
\usepackage{tikz-cd}
\input{translation-macros.tex}
\setcounter{secnumdepth}{3}
\begin{document}
\begin{center}{\Large Finite Element Exterior Calculus}\\代表性样本: Section 2.4--2.6 正文\end{center}
\noindent 原文定位: PDF 25--28, 印刷页 15--18. 本样本不包含之前的 Section 2.3 或之后独立的 Box 2.1. 文末仅提供已核实的样本外引用定位, 不表示那些章节已经翻译.\par
\setcounter{section}{1}
\section{同调代数的基本概念}\label{chap:homological-algebra}
\setcounter{subsection}{3}
% Agent-authored source transcription and faithful translation, verified against PDF 25-28.
% The sample starts at Section 2.4; it excludes the preceding Section 2.3 and later Box 2.1.
% block: chain-maps-heading page:25 type:section
\subsection{{链映射}}\label{sec:chain-maps}

% block: chain-map-definition page:25 type:ordinary-paragraph
{从链复形 $(V,\partial)$ 到另一个链复形 $(W,\xi)$ 的\emph{链映射}是一列线性映射 $f_k:V_k\to W_k$, 使得下图交换:}

% block: chain-map-diagram page:25 type:equation
\[
\begin{tikzcd}[column sep=1.2em,row sep=2.2em]
\cdots \arrow[r] & V_{k+1} \arrow[r,"\partial_{k+1}"] \arrow[d,"f_{k+1}"'] & V_k \arrow[r,"\partial_k"] \arrow[d,"f_k"'] & V_{k-1} \arrow[r] \arrow[d,"f_{k-1}"'] & \cdots \\
\cdots \arrow[r] & W_{k+1} \arrow[r,"\xi_{k+1}"'] & W_k \arrow[r,"\xi_k"'] & W_{k-1} \arrow[r] & \cdots
\end{tikzcd}
\]

% block: chain-map-induced-homology page:25 type:ordinary-paragraph
{换言之, 链映射是次数为 0 且与微分交换的分次线性映射, 即 $f\circ\partial=\xi\circ f$. 由于对 $c\in V_{k+1}$ 有 $f(\partial c)=\xi f(c)$, $f$ 将 $V$ 的 $k$-边界链映为 $W$ 的 $k$-边界链; 对 $c\in\mathcal Z_k$ 应用同一等式, 可知 $f$ 将 $k$-闭链映为 $k$-闭链. 因而链映射 $f$ 在同调上诱导线性映射 $\bar f$, 即映射 $\bar f_k:\mathcal H_k(V)\to\mathcal H_k(W)$, 它将某个 $v\in V_k$ 的同调类 $\bar v$ 映为 $f_kv\in W_k$ 的同调类 $\overline{f_kv}$. 这一映射具有函子性. 也就是说, 若 $(X,\zeta)$ 是第三个链复形, 且 $g$ 是从 $W$ 到 $X$ 的链映射, 则 $g\circ f$ 是从 $V$ 到 $X$ 的链映射, 并且 $\overline{g\circ f}=\bar g\circ\bar f$.}

% block: chain-projection-definition page:25 type:ordinary-paragraph
{后面我们将特别关注这样一种情形: 复形 $W$ 是 $V$ 的子复形, 且映射 $f_k$ 是 $V_k$ 到 $W_k$ 上的投影, 即对 $v\in W_k$ 有 $f_kv=v$. 此时称 $f$ 为\emph{链投影}.}

% block: chain-projection-surjectivity page:25 type:proposition
\begin{Proposition}\label{prop:chain-projection-homology-surjection}
{从复形到子复形的链投影在同调上诱导一个满射.}
\end{Proposition}

% block: chain-projection-surjectivity-proof page:25 type:proof
\begin{Proof}
{子复形的包含映射 $W_k\hookrightarrow V_k$ 显然构成一个链映射 $i$. 当 $f$ 是投影时, $f\circ i$ 是从 $W$ 到自身的恒等映射, 因而 $\bar f\circ\bar i=\overline{f\circ i}$ 是同调上的恒等映射, 特别地, 它是满射. 因此 $\bar f$ 必为满射.}
\end{Proof}

% block: cochain-complexes-heading page:26 type:section
\subsection{{上链复形}}\label{sec:cochain-complexes}

% block: cochain-complex-definition page:26 type:ordinary-paragraph
{上链复形与链复形相同, 只是指标的方向相反: 上链微分的次数是 $+1$ 而非 $-1$. 上链复形的指标通常写作上标而非下标, 并且常用 $d$ 而非 $\partial$ 表示微分:}

% block: cochain-complex-sequence page:26 type:equation
\[
\cdots\to V^{k-1}\xrightarrow{d^{k-1}}V^k\xrightarrow{d^k}V^{k+1}\to\cdots.
\]

% block: cohomology-definition page:26 type:ordinary-paragraph
{当然, 上述所有定义和结果都以显然的方式相应成立. 例如, 第 $k$ 个上同调空间 $\mathcal H^k$ 定义为闭上链 $\mathcal Z^k$ 模去上边界链 $\mathcal B^k$ 所得的商.}

% block: dual-cochain-complex page:26 type:ordinary-paragraph
{由于向量空间之间的任意线性映射 $L:V\to W$ 都在对偶空间之间诱导一个对偶映射 $L':W'\to V'$, 任意链复形都产生一个对偶上链复形. 例如, 从与单纯复形关联的单纯链复形出发, 我们得到单纯上链复形. 由于 $k$ 维单形构成 $C_k$ 的一组基, $k$-上链是 $C^k:=C'_k$ 的元素, 即一个函数 $\omega$, 它给每个 $k$ 维单形赋予一个实数 $\omega(f)$ (再由线性性扩张到链). 微分 $\partial^k:C^k\to C^{k+1}$ 由 $\partial^k=\partial'_{k+1}$ 定义, 即}

% block: dual-differential-evaluation page:26 type:equation
\[
(\partial^k\omega)g=\omega(\partial_{k+1}g),\qquad \omega\in C^k,\quad g\in C_{k+1}.
\]

% block: incidence-transpose page:26 type:ordinary-paragraph
{这意味着算子 $\partial^k$ 的矩阵是表示 $\partial_{k+1}$ 的定向关联矩阵的转置. 例如, 将 $\partial^0$ 作用于一个顶点, 所得的是与该顶点关联的各条边之和, 每条边根据其指向该顶点还是背离该顶点而分别乘以 $1$ 或 $-1$. 当然, 单纯上链复形}

% block: simplicial-cochain-complex page:26 type:equation
\begin{equation}\label{eq:simplicial-cochain-complex}
0\to C^0\xrightarrow{\partial^0}C^1\xrightarrow{\partial^1}\cdots\xrightarrow{\partial^{n-1}}C^n\to0
\end{equation}

% block: cohomology-duality page:26 type:ordinary-paragraph
{所包含的信息与单纯链复形所包含的信息本质上相同. 特别地, 可以证明上同调空间 $\mathcal H^k$ 自然同构于相应同调空间的对偶 $\mathcal H'_k$, 因而其维数是第 $k$ 个 Betti 数.}

% block: de-rham-heading page:26 type:section
\subsection{{Example: de Rham 复形}}\label{sec:de-rham-complex}

% block: de-rham-introduction page:26 type:ordinary-paragraph
{我们还将更详细地讨论 de Rham 复形}

% block: differential-forms-complex page:26 type:equation
\begin{equation}\label{eq:differential-forms-de-rham-complex}
0\to\Lambda^0(\Omega)\xrightarrow{d^0}\Lambda^1(\Omega)\xrightarrow{d^1}\cdots\xrightarrow{d^{n-1}}\Lambda^n(\Omega)\to0
\end{equation}

% block: de-rham-background page:26 type:ordinary-paragraph
{它与任意 $n$ 维区域或光滑流形 $\Omega$ 关联. 这里 $\Lambda^k(\Omega)$ 是 $\Omega$ 上的 (光滑) 微分 $k$-形式空间, 算子 $d$ 是由 Élie Cartan [\ref{bib:cartan-differential-expressions}] 引入的外微分. 我们将微分形式语言的介绍留到 Chapter~\ref{chap:exterior-calculus}. 目前只描述 $\Omega$ 是 $\mathbb R^3$ 中区域的情形, 在这种情形下, de Rham 复形可以用向量分析的术语来描述.}

% block: vector-calculus-complex-intro page:26 type:ordinary-paragraph
{此时该复形为}

% block: vector-calculus-de-rham-complex page:26 type:equation
\begin{equation}\label{eq:vector-calculus-de-rham-complex}
0\to C^\infty(\Omega)\xrightarrow{\operatorname{grad}}C^\infty(\Omega;\mathbb R^3)\xrightarrow{\operatorname{curl}}C^\infty(\Omega;\mathbb R^3)\xrightarrow{\operatorname{div}}C^\infty(\Omega)\to0.
\end{equation}

% block: calculus-identities page:27 type:ordinary-paragraph
{微积分恒等式 $\operatorname{curl}\operatorname{grad}\phi=0$ 和 $\operatorname{div}\operatorname{curl}u=0$ 保证 \eqref{eq:vector-calculus-de-rham-complex} 确实是一个复形. 它的上同调空间 $\mathcal H^k_{\mathrm{dR}}$ 构成该区域的 de Rham 上同调.}

% block: de-rham-map-zero-cochains page:27 type:ordinary-paragraph
{若区域 $\Omega$ 由某个单纯复形 $\mathcal S$ 给出三角剖分, 则存在一个从 de Rham 复形到单纯上链复形的有趣的上链映射, 称为 \emph{de Rham 映射}. 具体地, 我们给函数 $\phi\in C^\infty(\Omega)$ 配上一个 0-上链 $F^0\phi$, 它由对每个 $x\in\Delta_0(\mathcal S)$ 都有 $F^0\phi(x)=\phi(x)$ 确定. 正如我们将在 Chapter~\ref{chap:exterior-calculus} 看到的, 微分 $k$-形式可以在定向 $k$ 维流形上积分, 这里是最简单的 $k=0$ 情形. 由线性性将其扩张到任意 0-链, 便有 $F^0\phi(\sum c_ix_i)=\sum c_i\phi(x_i)$. 类似地, 我们给向量场 $\psi\in C^\infty(\Omega;\mathbb R^3)$ 配上一个 1-上链, 它在任意定向边 $e\in\Delta_1(\mathcal S)$ 上的值为}

% block: de-rham-edge-integral page:27 type:equation
\[
F^1\psi(e)=\int_e\psi\cdot t\,ds_e,\qquad e\in\Delta_1(\mathcal S),
\]

% block: de-rham-trace-forms page:27 type:ordinary-paragraph
{其中 $t$ 表示沿边 $e$ 的定向取向的单位切向量. 向量场 $\psi$ 可以看作微分 1-形式, 它在 $e$ 上有迹, 这个迹是一维流形 $e$ 上的微分 1-形式, 因而可以在其上积分. 上述线积分就是该积分用向量分析术语表示的形式. $\mathbb R^3$ 中的向量场也可以看作微分 2-形式, 此时它的迹在曲面上的积分有定义, 并可表示为}

% block: de-rham-face-integral page:27 type:equation
\[
F^2\psi(f)=\int_f\psi\cdot n\,ds_f,\qquad f\in\Delta_2(\mathcal S),
\]

% block: de-rham-volume-form-intro page:27 type:ordinary-paragraph
{这里 $n$ 是面 $f$ 的单位法向量, 其方向由 $f$ 的定向按照右手规则确定. 最后, 标量函数 $\phi$ 可以看作 3-形式, 从而得到}

% block: de-rham-volume-integral page:27 type:equation
\[
F^3\phi(T)=\int_T\phi\,dx
\]

% block: de-rham-connecting-map page:27 type:ordinary-paragraph
{其中 $T\in\Delta_3(\mathcal S)$ 是正定向四面体. 这样, 我们就定义了将 de Rham 复形 \eqref{eq:vector-calculus-de-rham-complex} 与单纯上链复形 \eqref{eq:simplicial-cochain-complex} 连接起来的 de Rham 映射:}

% block: de-rham-commuting-diagram page:27 type:equation
\begin{equation}\label{eq:de-rham-commuting-map}
\begin{tikzcd}[column sep=1.1em,row sep=2.2em]
0 \arrow[r] & C^\infty(\Omega) \arrow[r,"\operatorname{grad}"] \arrow[d,"F^0"'] & C^\infty(\Omega;\mathbb R^3) \arrow[r,"\operatorname{curl}"] \arrow[d,"F^1"'] & C^\infty(\Omega;\mathbb R^3) \arrow[r,"\operatorname{div}"] \arrow[d,"F^2"'] & C^\infty(\Omega) \arrow[r] \arrow[d,"F^3"'] & 0 \\
0 \arrow[r] & C^0 \arrow[r,"\partial^0"'] & C^1 \arrow[r,"\partial^1"'] & C^2 \arrow[r,"\partial^2"'] & C^3 \arrow[r] & 0
\end{tikzcd}
\end{equation}

% block: de-rham-first-cell-verification page:27 type:ordinary-paragraph
{de Rham 映射是上链映射, 也就是说, 该图交换. 例如, 我们来验证 $F^1\circ\operatorname{grad}=\partial^0\circ F^0$. 必须证明, 对任意光滑函数 $\phi$, 1-上链 $F^1\operatorname{grad}\phi$ 与 $\partial^0F^0\phi$ 相同, 即}

% block: de-rham-first-cell-identity page:27 type:equation
\[
F^1\operatorname{grad}\phi([x_0,x_1])=\partial^0F^0\phi([x_0,x_1])
\]

% block: oriented-edge-unit-tangent pages:27,28 type:ordinary-paragraph
{对该复形中的任意定向边 $[x_0,x_1]$ 都成立. 令 $t$ 表示单位向量 $(x_1-x_0)/|x_1-x_0|$, 则有}

% block: first-cell-calculation page:28 type:equation
\[
\begin{aligned}
F^1\operatorname{grad}\phi([x_0,x_1])
&=\int_{[x_0,x_1]}\operatorname{grad}\phi\cdot t
&&\text{{($F^1$ 的定义)}}\\
&=\phi(x_1)-\phi(x_0)
&&\text{{(微积分基本定理)}}\\
&=F^0\phi([x_1]-[x_0])
&&\text{{($F^0$ 的定义)}}\\
&=F^0\phi(\partial_1[x_0,x_1])
&&\text{{($\partial_1$ 的定义)}}\\
&=\partial^0F^0\phi([x_0,x_1])
&&(\partial^0=\partial'_1).
\end{aligned}
\]

% block: remaining-cells page:28 type:ordinary-paragraph
{图中另外两个方格的交换性可以用同样的方法证明, 其中用到向量分析中的 Stokes 定理和散度定理:}

% block: stokes-divergence-identities page:28 type:equation
\[
\int_f\operatorname{curl}\psi\cdot n=\int_{\partial f}\psi\cdot t,\qquad
\int_T\operatorname{div}\psi=\int_{\partial T}\psi\cdot n.
\]

% block: comparison-of-complexes page:28 type:ordinary-paragraph
{构成图 \eqref{eq:de-rham-commuting-map} 第一行的 de Rham 复形与构成第二行的单纯上链复形很不相同. 前者由通过微分算子连接的无限维函数空间组成, 而后者由有限维空间组成, 其算子本质上是组合性的 (乘以元素均为 $1$, $-1$ 或 $0$ 的矩阵). 然而, 两个复形的上同调相同. 这一陈述是 de Rham 定理的一种形式.}

% block: de-rham-isomorphism page:28 type:theorem
\begin{Theorem}{de Rham}\label{thm:de-rham-isomorphism}
{由 de Rham 映射在上同调上诱导的映射, 是从 de Rham 上同调空间到单纯上同调空间上的同构.}
\end{Theorem}

% block: singular-and-simplicial-cohomology page:28 type:remark
\begin{Remark}\label{rem:singular-simplicial-de-rham}
{刚才陈述的 De Rham 定理将借助微分算子用分析方法构造的 de Rham 上同调, 与由单形的链构造的、更具几何性的单纯上同调等同起来. 虽然单纯同调在某种意义上是最简单且最易计算的同调理论, 拓扑学家却往往更倾向于奇异同调, 因为它具有更大的普遍性和理论灵活性. 特别地, 单纯同调只对能够三角剖分的空间 (或这类空间的同胚像) 有定义; 而奇异同调由奇异单形 (单形的任意连续像) 的链构造而成, 对任意拓扑空间都有定义. 在大多数同调论的论述中, 一个较早出现且相当初等的结果是: 当两种方法都适用时, 即对经过三角剖分的空间, 它们给出相同的同调 (或上同调) 空间. 随后通常将更深刻的 De Rham 定理表述为: 在光滑流形上, de Rham 上同调与奇异上同调之间存在同构.}
\end{Remark}

% block: de-rham-later-proof page:28 type:ordinary-paragraph
{De Rham 定理是一个深刻的结果. 然而, 我们将在 Section~\ref{sec:whitney-forms-de-rham} 看到, 一旦备齐 FEEC 的工具, 就可以相当容易地证明它的一个版本.}

\clearpage
\section*{样本外引用定位}
\noindent 以下只登记原书引用目标, 并非这些部分的译文.\par
\subsection*{Chapter 6: Basic notions of exterior calculus}
\makeatletter\phantomsection\def\@currentlabel{6}\label{chap:exterior-calculus}\makeatother
\noindent Source PDF 73 / printed page 63.\par
\subsection*{Section 7.6: The Whitney forms $\mathcal P^-_1\Lambda^k$ and de Rham's theorem}
\makeatletter\phantomsection\def\@currentlabel{7.6}\label{sec:whitney-forms-de-rham}\makeatother
\noindent Source PDF 100 / printed page 90.\par
\subsection*{Reference [27]}
\makeatletter\phantomsection\def\@currentlabel{27}\label{bib:cartan-differential-expressions}\makeatother
\noindent Élie Cartan, Sur certaines expressions différentielles et le problème de Pfaff, Ann. Sci. École Norm. Sup. (3) 16 (1899), 239--332. MR 1508969. Source PDF 124 / printed page 114.\par
\end{document}
